% ============================================================
% 深思辨·悦读者（杂志页码 38–41）
% ============================================================
\colchapter{深思辨·悦读者}

\section*{迷茫不是台阶}
\auth{文/曲玮玮}

有个小朋友之前问了我一个问题，说有个作家来她学校做讲座，她提问说自己在迷茫期，不知道未来的出路在哪里，于是每天过得很懒散，感觉迷茫已经成了年轻人的通病，不知道自己该不该克服迷茫。其实，她提出这个问题的时候，心里已经有了预设答案，以为作家肯定会安慰她一番，甚至讲出“谁的青春不迷茫”这样的论调，然后让她趁着迷茫期多接触新事物之类。

提出这个问题，她只是为了迷茫得更心安理得。没想到作家劈头盖脸把她说了一顿：“什么叫迷茫是年轻人的通病？那些不迷茫的年轻人已经越跑越快了，你还在犹豫该不该克服迷茫？”然后，作家讲了自己的例子，说正是从小就找到了写作的使命，一步步努力才有今天的成就云云。

这个小朋友问我，迷茫不是好事吗，不迷茫的话怎么去探索人生啊？的确，像当年孟德尔实在不知道自己干吗，就闲着没事开始种豌豆，最后竟然种出了遗传定律。还有好多作家，也不知道自己适合干什么，就找张纸开始写文章，没想到真的写出了名堂。但迷茫始终不能成为逃避现实的借口。

年轻人特别容易沉湎于一种拧巴的迷茫，一方面觉得自己无所不能，另一方面觉得自己毫无长处。于是今天打打游戏，明天谈恋爱，后天去打个耳洞搞个文身，理所当然地浑浑噩噩，并自我安慰般想着——“如果年轻的时候就找到了自己的事业和使命，那一辈子就会走在一条单行线上，这样的人生也太无趣了。”毕竟，大器晚成的例子太多太多，所以自己也还早着。

但其实，那些“牛人”都很少迷茫，至少不会放纵自己永远处于迷茫的状态。太多科学家从小就对自己要探索的领域感兴趣，从而很早就为此付出大量时间。

在很多人还不知道自己想干吗的年纪，他们可能已经小有成就。想到冯唐也说，“写作是要有童子功的。”深以为然，有时候觉得自己若不是从小就开始读书写作，慢慢有了驾驭语言的能力，现在也不会这么从容。

很多时候你能成就一件事，是因为很久以前就埋下了种子。早点知道自己要做什么，人生省去了那么多时间的兜兜转转，就会清晰很多。当然，我也理解那些没有清晰目标感的人，在迷茫中去求索生活的意义，又何尝不是意义。迷茫本身没有错，错的只是你失去了走出来的斗志，用迷茫掩饰自己的惰性，甚至在这样尴尬的状态中沾沾自喜。

分享一段我很喜欢的话——人生就像一张长长的画卷，你很难一眼看到尽头。当下的努力很难找到一个明确的目的，却能让你变成更好的自己。既然迷茫，索性就不要东张西望，找找自己喜欢的事情做，练好本领，然后，等风来。

\begin{quoteBox}{洞见}
或许，迷茫是一种生命的必经，尤其在生活节奏越来越快、竞争压力越来越大的当下，谁的青春不迷茫？只是，有些人却把短暂的迷茫当作心安理得“躺平”的借口，渴望得到同情、安慰甚至施舍，他们安于懒惰的状态，很难真正认清自我并从挫折中重新站立。
\end{quoteBox}

\clearpage
\section*{自洽真的是人生万能解药吗？}
\auth{文/陶瓷兔子}

忘记了从什么时候开始，“自洽”被抬到一个煊赫的高度，很多的励志书籍甚至直接把“自洽”两个字当作万能解药，只要能自洽，你就可以获得终极的幸福。

我并不反对这两个字本身的意义——“我跟我相处得很和谐”，听上去好像没什么了不起，但和自己好好相处，哪怕只有一分钟其实已经相当不易。毕竟，人有太多欲望，又受太多鞭策，以至于对很多人来讲，跟自己相处就等于是超我和本我的互相折磨与互相指责。

不过，“与自我相处”绝不是人生的全部。试想有这样一个人，他的内在逻辑极其自洽，当他不想上进就不努力，考得很差也不会焦虑或觉得羞耻，没有工作也不为未来担心，有钱多吃点没钱少吃点，想睡觉就睡觉，想旅行就徒步前去，轻盈得好像一颗种子，随便被风吹到哪里都安之若素——很难找到比这更“自洽”的人了吧，可“我跟我相处得很和谐”的前提，是我与我的相处永远不会有第三方力量的介入——亲戚的议论，父母的叹息，同伴的压力，朋友的疏离，以及沉甸甸的整个社会评价体系。

没有多少人可以幸运地始终保持足够的轻盈去对抗沉重的现实，王小波那段话说得多好啊——“那一天我二十一岁，在我一生的黄金时代，我有好多奢望。我想爱，想吃，还想在一瞬间变成天上半明半暗的云。后来我才知道，生活就是一个缓慢受锤的过程，人一天天老下去，奢望也一天天消失，最后变得像挨了锤的牛一样。”我不知道在某种“试想”中很自洽的人，会不会在某个夜里因为留不住所爱之人，或者无力支付父母的养老开支被悔恨的暗影吞噬，总之在我有限的三十年人生里，我的的确确无数次因为虚度而懊悔不已——谁在浪费时间的当下不会觉得有种放松的欢愉呢？可人生之路就是得走到下一个阶段回头看，才能意识到哪些是铺路的梯，哪些是垫脚的石。

因此，“自洽”这个词不仅意味着此时此刻与自己相处得很和谐，还意味着在五年、十年之后回望这一天，依然只有欣慰而不会后悔。仅仅是这一点就太难了，而对于更多人来讲，一个藏得更深的难题是“不知道怎样才是真正的自我”。时间在叠加，自我会成长，这两个变量的不可控让“自洽”变成了一个伪命题，总有人可以自洽，但没有人可以一直自洽。

从心理学的角度来讲，我更倾向于用“稳态”这个词来代替“自洽”，像是走了很远很陡峭的山路终于来到一块平坦的栖息地，你可以在这里放松休整，但同时心知肚明那绝不是最终的目的地。主动也罢，被动也好，我们迟早会被自我或者生活推着离开那块平地，从稳态进入下一个坡道，过程中必然伴随着疲惫、迷茫和痛苦，而那也只是生活中最为正常的一部分而已。

\begin{quoteBox}{洞见}
绝对的自我，是要无条件遵从自己的内心，做到每时每刻的随心所欲，且完全地了解自己，而这无异于一种自欺欺人。其实，我们只要明白，世上绝没有一种全知全能的心理状态能让我们解决生活中的所有难题，就足以看清“自洽”这道伪命题。拥抱生活中的美与痛，接受它的全部，才是一个人成长并与自己相处的正道。
\end{quoteBox}

\clearpage
\section*{他，努力生长又决然远行}
\auth{文/吴淑斌}

陈昌雨控诉父亲纵火，是一桩旧案。今年以来，在陈昌雨的努力下，案件得到的关注越来越多，我们也才得以知晓：在云南曲靖的一个小山村里，一个父亲纵火点燃了房屋，致使妻子烧伤后治疗无效去世。随后，亲生儿子将父亲告上法庭，请求法院判处死刑，而这些举动也让他在亲戚当中颇受责备。

整件事情里最打动我的，正是陈昌雨本人。他就像一株从石缝里生长出来的小草，生活如此艰难，却还是在努力。我去到他租住的小屋子时，着实是被其间的干净整洁惊愕到：零食要整齐地码在架子上，吃完饭要立刻擦桌子拖地；他有两个空红酒瓶，插着绿植，小窗台上养了一盆多肉；房间里还有一个小鱼缸，里面有水草、鱼和增氧器。稿件发出后，有不少人托我给陈昌雨转点钱，他却拒绝了，“现在有工作……好意我心领了。”

一个原生家庭这么糟糕的孩子，为什么能成长得内心健康、强大？后来我想，大概正是因为父亲在他的童年里是缺失的。和母亲生活在一起，除了贫穷一些，他反而过得更松弛，得到了许多关爱。虽然再没有机会亲自见到他的母亲，但从陈昌雨身上，我相信他所描述的那个坚强、善良的母亲，确实如此。

只不过，在这对母子生活了十几年的村庄里，那些最靠近这家人的村民，态度却模糊不清。大概是案子慢慢在网上受到关注，为了维护乡村的颜面，人们似乎都在有意回避些什么。他们对陈昌雨“搞事情”“曝光家丑”的行为十分不满，认为他把父亲陈继卫告上法庭，更是“绝情”。另外，人们对陈继卫的评价倒是不错，即便他有两次入狱的经历，但那都不是发生在本村，好歹“他还给家里盖了房子”。

其实，这都是我意料之中的回应，一些再熟悉不过的陈词滥调。这部分村民依然在固守这种多年来留下的传统——男性的首要职责是为家庭挣得经济收入，女性则要维系家庭关系、照顾家人。当一个男人能够为家庭挣钱、盖房时，无论是否做好了丈夫、父亲的角色，他得到的评价都不至于太差劲。至于夫妻打架，则是再常见不过的“家务事”。这些村民当然不是“坏人”，相反，他们淳朴、好客，但当一件事情落入他们所熟悉的经验，能够用流传多年的传统来解释和处理时，很难要求他们“大义灭亲”，用冷冰冰的法律和规则来应对。

要和传统乡村决裂，需要巨大的勇气。坦白说，陈昌雨其实也没有做好这种准备。我的采访结束之后，他着急地打电话问我：“村里人怎么说？”他依然十分在意乡村社会的评价，在这里生活了十几年，他还无法熟练地用城市里学来的规则和秩序，去应对村庄生活。不过他依然是幸运的，工厂工作的机会，让他在与陈旧思想决裂的路上迈开了步子。

\begin{quoteBox}{洞见}
这几年，乡村仿佛一个承载着情感、记忆和经验的容器。人们在城市打拼，受够了“996”和“内卷”，像是有一种力量在号召我们离开城市、回到乡村。但在上述极端案例中，人们也骤然发现，村庄并不只有那些美好的场景，还有难以厘清的人情、捉摸不透的规则。从这个意义上说，渺小的个体很难与一套根深蒂固的生存法则相抗衡，但也总有人如陈昌雨一样以微弱的力量坚定地迈出独立思想的一步又一步。
\end{quoteBox}

\clearpage
\section*{延迟满足还是及时行乐？}
\auth{文/社长的小号}

“要不再等等吧，反正现在也不急着用”“熬过这一年，等上了大学就轻松了”“先将就着用，等以后赚钱了自己想买什么就买什么”……相信我们都（被）用过类似的话来克制自己的欲望，这些话的本质是教我们学会延迟满足。而与延迟满足相对的概念是及时行乐，这两种理念都有着各自的拥趸。

实际上，及时行乐和延迟满足并不冲突，当下的快乐也并不总以长远的幸福为代价，真正与延迟满足对立的，是被曲解和污名化的“及时行乐”。诸如“女人应该对自己好一点”等等消费话术就是生活中最常见的以“及时行乐”之名行不轨之事的典型，也是消解“及时行乐”本身积极意义的罪魁祸首之一。在这种消费主义的长期浸染下，人们不自觉地把及时行乐等同于即时满足物质欲望或是感官刺激，并将其视为实现长远目标道路上的绊脚石。但如果我们简单溯源一下就会发现，及时行乐是教人们把握当下，珍惜眼前，而非不计后果地自我放纵。

那么，到底是等枝头挂满硕果，再享受收获的喜悦；还是在蓓蕾初绽时，就赶紧摘下一饱眼福？不做选择的成年人可以两者兼得吗？

尽管我们经常将延迟满足和自律挂钩，但是研究发现，自我控制以及由此带来的所有好处，可能与抑制冲动完全无关。因为，与想象中不同，掌握着“自律”话语权的人们谈到自己的好习惯时，用的并不是坚持、忍受这些词汇。相反，他们更乐意用上瘾、享受来解释自己的十年如一日。在从事某种活动的过程中，他们品尝到了满足感，这种满足感会激励他们延续行动，并由此形成良性循环。所以，我们不妨将自律神话丢在一边，然后跟随哲学家爱比克泰德的指引，思考一下“我是谁”“将要到哪去”这两个“终极问题”，并根据得到的答案，去重构自己的快乐体系。

另外，延迟满足不可避免会带来痛苦，但是痛苦自有其存在的价值。功利主义伦理学家穆勒认为，“成为一个不高兴的人总好过一头高兴的猪。”快乐和痛苦是一对双生花，正是旷日持久的参与和坚持不懈的付出这些看似痛苦的行为，使得快乐的感觉更加深刻而持久。高质量的快乐是从无到有，通过自己参与和创造而获得的，在这个过程中我们所经历的痛苦恰恰是我们存在和实践的证明。

最后，快乐是有时效性的，就像小时候喜欢但没有得到的玩具，长大后即便可以轻易得到，却已经失去了当初的快乐。所以，一味推崇“延迟满足”，不仅会产生逆反心理，让时间成本沉没，还可能导致在过程中感到挫败、变得越来越优柔寡断、更加贪婪……在一眼望不到头的道路上奔走，有时候难免需要世俗的、可以立即满足我们的快乐作为慰藉，这时候也不必过于苛责和压抑自己的欲望。难过的时候吃一块平时不敢吃的蛋糕，为一部追更了很久的剧熬一次夜，谁说享受当下的快乐不会成为我们重新出发的动力呢？

\begin{quoteBox}{洞见}
我们既需要当下简单的快乐，也追求持续付出后得以成功的满足；想旅行时我们可以说走就走，而对事业的热爱同样是我们穷尽一生的习惯……压抑欲望不代表终将厚积薄发，即时满足也不意味着无度挥霍、沉溺放纵——对自己负责，永远不拘于形式，而是量体裁衣、量力而行的顺其自然。
\end{quoteBox}
